\begin{abstract}
We report a bounded mechanistic case study of a compositional generalization failure
in T5-small fine-tuned on SCAN \texttt{add\_prim\_jump}.
At the checkpoint studied, the model fails on 56.8\% of compound-jump test commands
despite the encoder marking \texttt{jump} as categorically out-of-distribution
($p=0.0000$).
The dominant failure mode is substituted-walk: the model generates the correct cycle
structure but substitutes \texttt{I\_WALK} for \texttt{I\_JUMP} at every action slot.
Decoder cross-attention routing to the relevant encoder position shows no detected
difference between the stored fail and success comparison groups ($p=0.9208$),
indicating that routing is not the primary discriminator in this comparison.
The failure is localized to the value projection.
Value-direction evidence and subsequent repair experiments identify decoder heads L4H6
and L5H5 as repair-relevant pro-walk value-projection sites.
Logit decomposition shows near-cancellation: a cross-attention total of $-171.25$
logit units is opposed by $+165.25$ units of FFN/residual correction, leaving a net
margin of $-6.005$ ($28.52\times$ reconstruction fraction).
Output embedding norm asymmetry ($\|\texttt{I\_WALK}\|\approx 520$ versus
$\|\texttt{I\_JUMP}\|\approx 424$) amplifies this small pro-walk residual to
$P(\texttt{I\_WALK})=0.914$.
A targeted rank-1 modification to $W_V$ at L4H6 and L5H5 shifts the sampled
fail-group margin by $+2.902$ logit units and flips 16 of 30 sampled failures,
using perturbations of approximately 8\% of each edited head's $W_V$ Frobenius norm,
with monotonic dose-response.
Controlled toy geometry (G-track) independently demonstrates near-cancellation
divergence and value-direction substitution; these results provide structural analogy,
not calibrated prediction.
All results are from one checkpoint on one SCAN split; the repair is partial; and
broad specificity against unrelated commands is not established.
\end{abstract}
